\documentclass[11pt,letterpaper]{article}
\usepackage[margin=1in]{geometry}
\usepackage[T1]{fontenc}
\usepackage{amsmath,amssymb,microtype}
\usepackage[hidelinks]{hyperref}
\pagestyle{empty}
\setlength{\parindent}{0pt}
\setlength{\parskip}{0.6em}
\hypersetup{pdftitle={Cover letter to Mathematics of Computation},pdfauthor={Haotian Zhong}}
\begin{document}
\textbf{Haotian Zhong}\\
Department of Mathematics, School of Sciences\\
Shanghai University, Shanghai 200444, P. R. China\\
\href{mailto:zhonghaotian1219@gmail.com}{zhonghaotian1219@gmail.com}\\
ORCID: 0009-0009-5567-1315

October 1, 2026

Editors of \textit{Mathematics of Computation}\\
American Mathematical Society

Dear Editors,

I submit the manuscript \textbf{``Frequency Saturation and Certified Computation of Relaxation Functionals''} for consideration in \textit{Mathematics of Computation}.

The paper concerns the computation of linear functionals of a positive distribution of relaxation times from impedance measurements, the Stieltjes-type inverse problem $Z(\omega)=R+i\omega L+\int(1+i\omega\tau)^{-1}\,d\mu(\tau)$ with unknown nonnegative resistance and inductance. It determines how much frequency information such a computation requires, and it gives a finite numerical method, with a complete termination proof, that computes the entire data-consistent range of the functional with certified numerical error.

The main results are the following. (1)~Exactly $K+2$ distinct frequencies are necessary and sufficient for uniform noiseless identification of references with at most $K$ atoms among all finite positive measures, and $K+1$ when the calibration parameters are known. (2)~At this threshold a fixed design attains the worst-case accuracy order of the entire positive-frequency response; the sharp exponent is determined by local target regularity and cluster multiplicity, uniformly through colliding atoms and vanishing masses. Positive moment-matched spectra that stay close at every frequency give the lower bound, including against target-aware adaptive queries, and calibration-annihilating polynomial reproduction gives the upper bound. (3)~An exact-arithmetic refinement algorithm encloses both endpoints of the continuous feasible range with any prescribed numerical excess; finite optimization gaps, target approximation, and continuum correction enter separate error budgets, and the certificates yield an a posteriori absolute--relative stopping rule. (4)~The same modulus gives matching expected acquisition costs for prescribed terminal width and confidence under a Gaussian variance--effort law, against adaptive frequency, precision, and stopping; a class-tuned expanding design that invokes the certified computation attains this order and terminates at every finite positive spectrum.

The work joins the analysis of a restricted information class to an executable, accuracy-controlled inverse computation, and so addresses approximation of positive measures, information complexity, and verified numerical methods. The numerical section reports exact full-frequency certificates, equal-effort design comparisons, certified cost bounds, and executed stopping paths; a reproducibility archive with exact inputs, verification code, and certificates accompanies the submission.

The manuscript is original, unpublished, and not under consideration elsewhere. I declare no conflicts of interest. I suggest Professor Markus Bachmayr as a handling editor, given his expertise in approximation, adaptive methods, and inverse problems.

Sincerely,\\[0.35em]
Haotian Zhong
\end{document}
